\section*{Abstract}

The rise of cryptocurrencies has been paralleled by an increase in fraudulent activities, posing significant challenges for financial ecosystems and regulatory bodies. This study explores fraud detection within cryptocurrency transactions using price data and social media activity as primary data sources. The research leverages open and close price metrics alongside embeddings extracted from posts and comments on Reddit and Twitter to construct predictive models. MultiRocket was successfully implemented as the primary machine learning model, demonstrating promise in detecting fraudulent patterns. Attempts to integrate LSTM models and the ConvTran architecture (via ConvTran) were initiated but remained incomplete, highlighting potential avenues for future work. The findings and methodologies presented here aim to establish a foundation for advancing fraud detection in cryptocurrency ecosystems, encouraging further development and innovation in the field.

\vspace{2ex}

\textbf{Keywords:}

Cryptocurrencies, Bitcoin, Fraud, Web Data Acquisition, Social Media, Reddit, X (Twitter), Time Series Analysis, MultiRocket, LSTM

\clearpage

\section*{Acknowledgment}

We, the authors, would like to extend our special thanks to Gabriel Torres Gamez for his collaboration with us on gathering and conducting exploratory data analysis (EDA) on social media data, as well as to Moritz Kirschmann and Adrian Brändli for their guidance throughout this Challenge-X.

\clearpage

\section*{Note}

The texts written in this paper may have been partially generated or edited with the assistance of ChatGPT \parencite{noauthor_chatgpt_2024}.
